\source{Xin et al., Findings of ACL 2026，PDF p.1}
\section{MetaMem}\presenter{YaoYang-123456}
\begin{frame}[plain]
\vspace{7mm}{\small\color{accent}专题论文阅读}\par
\vspace{7mm}{\fontsize{28}{34}\selectfont\bfseries MetaMem}\par
\vspace{4mm}{\Large 让模型正确使用已取回的事实}\par
\vspace{7mm}\textcolor{accent}{\rule{\textwidth}{0.8pt}}\par
\vspace{3mm}YaoYang-123456\hfill 约8分钟核心讲解
\note{大家好，我是第八队的YaoYang-123456。这一部分介绍 MetaMem，发表于二零二六年 ACL Findings。它研究的是长期对话问答：模型已经取回相关事实，为什么仍然回答错误？作者在具体事实之外，再学习一组指导事实使用的经验。下面先用一个计数例子说明问题，再解释规则怎样产生，以及实验能够支持什么结论。}
\end{frame}

\source{原文 PDF p.1，图1；图中里程数字按论文示例整理}
\begin{frame}{检索到了事实，为什么还会答错？}
\small
\textbf{问题：四次旅行共走了多少英里？}\par\vspace{3mm}
\begin{tabularx}{\linewidth}{@{}p{34mm}Y@{}}
\toprule
事实 A & 前三次旅行合计\key{1800}英里\\[1mm]
事实 B & 第四次旅行总计\key{1200}英里\\[1mm]
事实 C & 第四次旅行第一天走了\key{300}英里\\\bottomrule
\end{tabularx}\par\vspace{4mm}
正确：$1800+1200=\mathbf{3000}$\quad 错误：$1800+1200+300=3300$\par
\aside{错误未必来自检索：事实可能全部存在，模型却把“总量”和“部分”重复相加。}
\note{请看页面上的旅行例子。用户问，四次公路旅行总共走了多少英里。取回的事实有三条：前三次旅行合计一千八百英里；第四次旅行总计一千二百英里；第四次第一天走了三百英里。这里最重要的关系是，第一天属于第四次旅行，因此三百已经包含在一千二百里面。把三个数字都相加会得到三千三百，正确做法是只加两个不重叠的总量，得到三千。这个例子来自论文，页面做了简化整理。它说明，检索找到了相关数字，并不保证模型理解这些数字的事件和包含关系。MetaMem就是从这一步的错误入手。}
\end{frame}

\source{原文 PDF p.3，§3.1，式(1)--(3)}
\begin{frame}{事实记忆 $M$ 与元记忆 $E$}
\small
\begin{tabularx}{\linewidth}{@{}p{23mm}YY@{}}
\toprule
对象 & 保存什么 & 例子\\\midrule
事实记忆 $M$ & 具体经历、数字、时间、人物 & 第四次旅行总计1200英里\\
元记忆 $E$ & 可迁移的知识使用策略 & 总量与部分不可重复计数\\\bottomrule
\end{tabularx}\par\vspace{5mm}
\[\hat y=\mathrm{LLM}(\operatorname{Prompt}(q,M,E))\]
\vspace{2mm}
\textbf{训练对象：}自然语言规则集合 $E$；不通过反向传播更新基础模型权重。\par
\aside{信息不足时，规则也不能补出原本没有检索到的关键事实。}
\note{论文把两种作用不同的记忆分开。事实记忆用字母 M 表示，记录发生过什么，比如某次旅行的时间和里程。元记忆用字母 E 表示，记录怎样使用事实的经验，例如先检查总量与部分是否重叠，再进行求和。这条经验不包含本题的具体答案，可以用于新的计数问题。请看公式：q 是当前问题，Prompt 表示把问题、事实和经验组织成提示词，LLM 是大语言模型，输出的 y 是模型回答。普通记忆问答主要输入问题和事实，MetaMem还输入学到的经验。这里学习的主要对象是外部自然语言规则，模型参数保持不变；但规则会改变模型接收到的上下文。可以把 M 看成资料，把 E 看成使用资料时的注意事项。注意事项也可能错误，必须根据反馈检查；它既不是额外的标准答案，也不是模型新训练出来的一组权重。这个区分决定了后面的更新操作为什么是编辑文字，而不是计算梯度。}
\end{frame}

\source{原文 PDF pp.3--5，Figure 2，式(4)--(10)}
\begin{frame}{从候选回答到反思与规则更新}
\small
\begin{enumerate}\setlength{\itemsep}{2mm}
\item \textbf{多次回答：}在同一事实记忆下采样若干候选回答。
\item \textbf{评价与反思：}使用训练标准答案判断对错并归纳原因。
\item \textbf{编辑规则：}提出 ADD / DEL / MOD，处理重复与冲突。
\item \textbf{验证与冻结：}用验证集选取规则版本；测试时不再更新。
\end{enumerate}
\vspace{2mm}
\begin{center}
\key{回答} $\rightarrow$ \key{反馈} $\rightarrow$ \key{反思} $\rightarrow$ \key{规则} $\rightarrow$ \key{新回答}
\end{center}
\aside{训练时可以看到标准答案；测试回答时不能看到答案，更不能用测试答案改规则。}
\note{接下来沿着一次训练看完整流程。先固定问题、事实和当前规则，生成若干候选回答。论文默认每题采样五次，可能同时出现正确的三千和错误的三千三百。评价模型叫 judge，它结合问题、标准答案、事实和回答给出正误反馈。反思模型再解释成功或失败的原因，例如错误回答把已包含的第一天里程重复计入。随后，模型根据这些反思提出规则编辑。一个批次可能有重复或者冲突的提案，因此先筛选，再执行，得到下一版元记忆。训练标准答案可以参与这些反馈；测试时只能使用提前选定并冻结的规则，不能用测试答案继续修改它。}
\end{frame}


\source{原文 PDF pp.4--5，Figure 2、式(4)--(10)；以下规则文字为教学构造}
\begin{frame}{一次自然语言规则编辑}
\small
\textbf{原规则：}把所有出现的数字相加。\par\vspace{2mm}
\textbf{反思：}第一天属于第四次旅行，重复计入会多算。\par\vspace{2mm}
\textbf{MOD 提案：}求和前核对对象、事件和包含关系，只合并互不重复的数量。\par\vspace{3mm}
\begin{tabularx}{\linewidth}{@{}p{23mm}Y@{}}
\toprule
ADD & 新增缺失的经验；从空集合开始时先添加\\
DEL & 删除错误或过时的经验\\
MOD & 修改已有编号的经验\\\bottomrule
\end{tabularx}
\aside{提案先处理重复和冲突，再执行到 $E$；例子不代表本组真实学出的规则。}
\note{这一页把符号优化展开成具体文字。假设已有规则写着，把所有出现的数字相加。反思指出，这句话没有区分总量与部分。模型可以提出修改，把它改成：求和前先核对对象、事件和包含关系，只合并互不重复的数量。这里的修改叫 MOD；新增经验叫 ADD；删除错误经验叫 DEL。假如另一条提案还要求保留所有分项逐项相加，就会与新规则冲突，需要在批量过滤时处理。最后执行的是筛选后的动作，而不是无条件把全部建议追加进去。这是为解释流程构造的例子，不是本组真实学出的规则。新问题出现时，模型仍需根据那道题的实际事实作答，不能直接套用旧数字。}
\end{frame}

\source{原文 PDF pp.5、11，§4，附录 A.3、A.5}
\begin{frame}{实验设置与数据隔离}
\small
\begin{tabularx}{\linewidth}{@{}p{34mm}Y@{}}
\toprule
数据和划分 & LongMemEval 500题；每折350训练＋50验证＋100测试，共五折\\[1mm]
事实记忆 & LightMem；规则以外还存在检索模块\\[1mm]
生成模型 & Qwen3-30B-A3B-Instruct、Llama3.1-70B-Instruct\\[1mm]
评价模型 & Qwen3-235B-A22B；判分是实验的一部分\\[1mm]
学习设置 & 5 epochs，每批50题，每题采样5次\\\bottomrule
\end{tabularx}
\aside{模型、数据划分、评价标准不一致时，不能直接把准确率差异归因到元记忆。}
\note{理解结果之前，先看实验条件。主要基准是 LongMemEval，可以理解为长期对话记忆问答，共五百题。论文采用五折：每折三百五十题用于学习规则，五十题验证，一百题测试。验证用于选择规则版本，测试用于评价选定版本。事实记忆主要用已有记忆系统 LightMem 构建，每次检索二十条；两种回答模型是表中的 Qwen 三三十 B 和 Llama 三点一七十 B，评价模型是更大的 Qwen 三。B 表示十亿参数，完整型号保留在表中。学习五轮，每批五十题、每题五次采样。五折是更换测试子集，五轮是在一折训练数据上重复学习，这两个数字不能混淆。准确率是被判正确的问题比例。正文把训练和验证合称四百题构建数据，附录才进一步拆开。}
\end{frame}

\source{原文 PDF p.6 Table 1、p.7 Table 2；均为论文结果}
\begin{frame}{主结果与消融的解释}
\small
\textbf{原论文 LongMemEval 平均值（\%）}\par\vspace{2mm}
\begin{tabularx}{\linewidth}{@{}Yrrr@{}}
\toprule
生成模型 & LightMem & MetaMem & 提升（点）\\\midrule
Qwen3-30B-A3B & 67.50 & \key{71.90} & +4.40\\
Llama3.1-70B & 66.17 & \key{69.08} & +2.91\\\bottomrule
\end{tabularx}\par\vspace{4mm}
\textbf{Qwen 消融：}完整方法 71.90；无反思 70.02；无演化 68.33。\par\vspace{3mm}
\aside{两组模型的增幅不同，整体提升不等于所有题型改善；这些数字不是本组试点成绩。}
\note{表中是论文结果，不是我们的复现成绩。相同 Qwen 设置下，LightMem 的平均值为百分之六十七点五零，加入 MetaMem 后为七十一点九零，增加四点四零个百分点。Llama 从六十六点一七到六十九点零八，增加二点九一个百分点。这里不能把百分点说成相对增长率，也不能说每个题型都提高了。接着看消融：Qwen 的完整方法为七十一点九零；去掉反思后是七十点零二；去掉演化后是六十八点三三。去掉反思的版本仍编辑规则，只是省略中间的原因分析。这支持两个环节在论文设置下有作用，但没有证明任意任务、任意小模型都能获得同样收益。}
\end{frame}

\source{原文 PDF §Limitations、附录 A.6；本组方案 docs/project-plan.md}
\begin{frame}{局限、适用条件与 StructMem 的联系}
\small
\begin{columns}[T,onlytextwidth]
\column{.47\textwidth}
\textbf{优势与价值}\par\vspace{2mm}
\begin{itemize}\setlength{\itemsep}{2mm}
\item 分离“检索缺失”和“利用错误”
\item 规则可读、可修改，无须微调权重
\item 消融实验提供一定证据
\end{itemize}
\column{.48\textwidth}
\textbf{限制与风险}\par\vspace{2mm}
\begin{itemize}\setlength{\itemsep}{2mm}
\item 依赖带答案数据与评判质量
\item 多次采样和反思有成本
\item 过度泛化、冲突与事实缺失
\end{itemize}
\end{columns}
\vspace{3mm}
\aside{StructMem 侧重组织和检索证据；MetaMem 侧重检索后的事实利用。两者不可直接混算实验成绩。}
\note{最后总结方法的价值和边界。优点是把取回事实和利用事实分开，并把使用经验写成可检查、可更新的规则。局限是学习仍依赖带答案样本和评价质量，多次采样、判分、反思与编辑都有成本；规则也可能冲突或过度泛化。关键事实没有找到时，规则无法凭空补齐证据。对本组而言，StructMem主要研究事实和关系如何组织、检索，MetaMem研究检索之后如何正确利用。个人简化试点可以帮助定位这两类错误，但不能与团队 LoCoMo 成绩混算，也没有证明规则带来稳定提升。我认为这篇论文最值得借鉴的是，把错误落实到证据与使用过程，再用受控实验检验改进。}
\end{frame}
